\section{Auxiliary results for Section~\ref{s:encoGames}}
%Operational correspondence between OGS and \piI
\label{a:oc}


We present results that are needed to establish the operational correspondence between 
OGS and \piI, studied in Section~\ref{s:encoGames}. 


We begin discussing an optimisation
$\qencV$ of
the  encoding of 
 call-by-value $\lambda$-calculus.
Following 
 Durier et al.'s 
\cite{DurierHS18}, 
we sometimes use $\qencV$  to simplify proofs. 
We report the 
full definition of   $\qencV$;
for  convenience, we also list the 
 definition of $\qencV$  on 
 OGS environments and configurations, thought  is the same as 
that for the initial encoding $\cbvSymb$~--- just replacing $\cbvSymb$ with 
 $\qencV$ in Figure~\ref{f:OGS_pi_cbv}). 


The encoding 
 $\qencV$
 is 
obtained  from the  initial one 
$\cbvSymb$ 
by inlining the encoding and performing a few (deterministic) reductions,
at the price of 
 a more complex definition.  Precisely, in the encoding of application 
 some of the initial
communications are removed, including those with which a term signals to have become a value. 
To achieve this,  the encoding of an 
application  goes by a case analysis
on the occurrences of values in the subterms.

% The general idea of the optimized encoding can be illustrated on two
% particular cases. In $\encoVV{\val M}$,  we essentially 
%  unfolding the original encoding, and then  performing one
% (deterministic) communication. 
% %
% In the case of $\encoVV{x\val}$, besides  unfolding the original
% encoding and performing some  deterministic communications, we also perform 
% the administrative reductions that always precede the execution of $\enca{x\val}$.
% \finish{in the figure, probably remove the encoding of environments and configurations} 

\begin{figure*}
\Mybar 
{\small
Encoding of  $\lambda$-terms:

\[
\setlength{\arraycolsep}{3pt}%
\begin{array}{rcl@{\;}l}
\equaDS{}
{  \encoV {x\val}} \defi{ \bind p \outb x {z,q}.(\encoVVa \val z|\fwd q p)}
\\
\equaDS{}
{ \encoV {(\lambda  x.  M)\val}} \defi{ \bind p \new {y,w} (\encoVVa {\lambda  x. M} y | \encoVVa {\val} w
  | \outb y {w',r'}.(\fwd {w'} w|\fwd {r'} p))} 
\\
\equaDS{{(*)}}
{ \encoV {\val M}} \defi{ \bind p \new y (\encoVVa \val y | \new r (\encoVa M r | \inp r w
  .\outb y {w',r'}.(\fwd {w'} w|\fwd {r'} p)))}
\\
\equaDS{{(*)}}
{ \encoV {M\val} } \defi {\bind p \new q (\encoVa M q | \inp q y. \new w (\encoVVa \val w |
  \outb y {w',r'}.(\fwd {w'} w|\fwd {r'} p)))}
\\
\equaDS{{(**)}}
{ \encoV {MN}} \defi {\bind p \new q (\encoVa M q | \inp q y. \new r (\encoVa N r | \inp r w
  .\outb y {w',r'}.(\fwd {w'} w|\fwd {r'} p)))}
\\
\equaDS{}
{ \encoV {\val}}\defi{ \bind p \outb p y . \encoVVa\val y}
\end{array} 
  \]
where
in the  rules marked $(*)$, $M$ is  not   a value,  and  
in the rule marked $(**)$   $M$ and $N$  are  not values;
and where 
 $ \encoVV\val$ is thus defined:
\[
\begin{array}{rcl}
\encoVV {\lambda  x. M}\defi{ \bind y  !\inp y {x,q}.\encoVa M q}\\
\encoVV x \defi {\bind y \fwd y x}
\end{array}
\]
$ $ \dotfill $ $ 


Encoding of environments: 
\[ 
\begin{array}{rcl}
\encoEnvV{\pmap{y}{V}\cdot\gamma'  } & \defi & \encoVVa {V} y | \encoEnvV{\gamma'  }
\\[\mypt] 
\encoEnvV{\pmap{q}{(E,p)}\cdot\gamma'  } & \defi & 
\inp qx .\encoVa {E[x]} p  | \encoEnvV{\gamma'  } \\[\mypt]
\encoEnvV{\emptymap  } & \defi & \nil
\end{array} 
\]
$ $ \dotfill $ $ 

Encoding of configurations:


\[ 
%Encoding of active configurations $\conf{M,p,\gamma}$: 
\begin{array}{rcl}
\encoConV{ \conf{M,p,\gamma}}
& \defi& 
 \pairP {
 \encoVa {M} p | \encoEnvV \gamma }
\\[\mypt]
%Encoding of passive configurations $\conf{\gamma}$: 
 \encoConV{ \conf{\gamma}} & \defi&  
 \pairP {
\encoEnvV \gamma}
\\[\mypt]
% Encoding of the initial configuration:
 \encoConV{ \conf{M }}
& \defi&  \pairP { \encoV   M }
\end{array}
 \] 
}
\caption{The optimised  encoding}
\label{f:opt_encod}
\Mybar 
\end{figure*}
 


\DSOCTb


We begin with  a few  results from \cite{DurierHS18,Durier20}
that are needed to reason about the  
optimised encoding $\qencV$, in addition to Lemma~\ref{l:opt_sound}
reported in the main text. 
%(e.g., Theorems~\ref{t:opcorrCON} and ~\ref{t:opcorrCON2}). 

\DSOCTe


\begin{lemma} 
\label{l:opt_stuck}
We have: 
$$
\begin{array}{rcl}
\encoVa {\evctxt[x\val]} p\exn \outb x {z,q}.(\encoVVa \val z|\inp q
y.\encoVa {\evctxt[y]}p).
\end{array}
 $$ 
\end{lemma}
\noindent
The proof \cite{DurierHS18,Durier20}
goes by induction on the evaluation context \evctxt.



The following lemma \cite{DurierHS18,Durier20}
uses Lemma~\ref{l:opt_stuck}  
to establish the shape of the possible transitions that a term 
$\encoVa M p $ can perform. 


\begin{lemma}
\label{l:opt_op}
For any
  $M\in\Lao$ and  $p$, process  $\encoVa M p$ has exactly
  one immediate transition and exactly one of the following clauses holds:
\begin{enumerate}
\item $\encoVa M p\arr{\outb py}P$  %for some $P$
and $M$ is a value, with $P=\encoVVa M y$; %(and therefore either $M=x$ for some variable $x$ or $M=\lambda y.N$ for some $y,N$), 
%and $P=\encoVVa \val y$
\item $\encoVa M p\arr{\outb x {z,q}} P$  and %for some $x$ and $P$, 
 $M $ is of the form $\evctxt[x\val]$, for some $E,x,\val$,  with 
$$P\exn \encoVVa \val z|\inp q y.\encoVa {\evctxt[y]} p ,$$
and moreover $z$ is not free in $\inp q y.\encoVa {\evctxt[y]} p$
whereas $q$  is not free in $\encoVVa \val z$;
  %for some value $\val$ and some evaluation context $\evctxt$%, and $\encoVa M p\sim\outb x {z,q}.P $
\item $\encoVa M p\arr\tau P$ and there is $N$  with $M\red N$ and
  $P\exn \encoVa N p$. 
\end{enumerate}
\end{lemma}
 
\noindent
We now move to reasoning about the behaviour of the encoding of configurations. 
Two key lemmas are the following ones; they are derived from Lemmas~\ref{l:opt_stuck} and
\ref{l:opt_op}.

\begin{lemma} %[Actions from $\pi$-calculus configurations]
\label{l:opt_con}
If  $\encoConV{\gamma } \arr \mu P$ and $\mu$ is an input action, then we have three possibilities:
\begin{enumerate}

\item $\mu $ is an input $y(x,q)$ and $\gamma 
= \pmap{y}{\lambda x . M}\cdot\gamma' $, and $P \contr 
\encoVa {M} q  | \encoEnvV{\gamma  } $; 

\item $\mu $ is an input $y(x,q)$ and $\gamma 
= \pmap{y}{z}\cdot\gamma' $, and $P \contr 
\encoVa {zx} q  | \encoEnvV{\gamma  } $; 


\item $\mu $ is an input $q(x)$ and $\gamma 
= \pmap{q}{(E,p)}\cdot\gamma' $, and $P\contr 
\encoVa {E[x]} p  | \encoEnvV{\gamma'  } $. 
\end{enumerate} 
\end{lemma}
 
% We will also use the following result, on the actions from a value. 
% \iflong
% \finish{This is a new result}
% \fi

\begin{lemma}
\label{l:VAL}
  $\encoVVa {\val} y  \arr {\inp y {x,q}} \wbPI 
\encoVa {\val x } q | \encoVVa {\val} y$
\end{lemma}              
\noindent
In Lemma~\ref{l:VAL}, the occurrence of $ \wbPI$ cannot be replaced by $\contr$; 
for instance, if $V = \lambda x. M$ then, using Lemma~\ref{l:opt_op}, we can infer that 
the derivative process is in the relation $\contr$ with 
$\encoVa {M } q$; however it is not in the same relation with 
$\encoVa {\val x } q$ (which has extra $\tau$ transitions with respect to $\encoVa {M } q$).

 
\iflong
\finish{below i think it is actually not needed now} 
Here is another important result that we borrow from \cite{DurierHS18}.
It shows  that the only $\lambda$-terms
 whose encoding is bisimilar to 
$\enca x$ reduce either to $x$, or to a (possibly infinite) 
$\eta$-expansion of $x$.

\begin{lemma}\label{l:opt_eta}
  If $\val$ is a value and $x$ a variable,
  $\encoVV \val\wbPI \encoVV x$ %if and only if
  implies that either $\val=x$ or $\val = \lambda   z. {M}$, where the eager
  normal form of $M$ is of the form $\evctxt[x\valp]$, with
  $\encoVV {\valp}\wbPI \encoVV z$ and $\encoV {\evctxt[y]}\wbPI\encoV y$
  for any $y$ fresh.
\end{lemma}

\fi



% %\paragraph{Operational correspondence.}
% \section{Auxiliary results for Section~\ref{s:encoGames}}
% \label{a:aux_tr_bisi_ALT}

\noindent
We can now establish the operational correspondence between OGS and \piI. 
% The results below are needed in the proofs of Theorems~\ref{t:opcorrCON_w} and
% \ref{t:opcorrCON2_w}, together with the Lemma~\ref{l:opt_sound}.
(In all the results about operational correspondence, we exploit the convention on
freshness of bound names of actions and traces produced by OGS configurations and \piI
processes, as by Remark~\ref{r:bn}.)
\begin{lemma}[strong transitions, from \AOGS to \piI]
\label{l:opcorrCON_str}
$ $ \begin{enumerate}

\item If $\FF \longrightarrowA \FF'$, then 
$ \encoConV \FF \longrightarrowPI  {\contr} \: \encoConV {\FF'}$;


\item If $\FF \arrA\ell \FF'$, then
$ \encoConV \FF \arrPI\ell  
 {\wbPI} \:
\encoConV {\FF'}$.
\end{enumerate} 
\end{lemma} 

\begin{proof}
The proof goes by a case analysis on the rule used to derive the transition from 
$\FF $, using the definition of the encoding, using Lemmas~\ref{l:opt_op}-\ref{l:VAL}.
\end{proof} 



\begin{lemma}[strong transitions,  from  \piI to \AOGS]
\label{l:opcorrCON2}
$ $
\begin{enumerate}
\item If 
$ \encoConV \FF \longrightarrowPI \PP$ then there is $\FF'$ such that 
$\FF \longrightarrowA \FF'$  and  
$\PP   \contr \encoConV {\FF'}$; 

\item If
$ \encoConV \FF \arrPI\mu \PP$
and $\mu$ is an output, 
then there is $\FF'$ such that
$\FF \arrA\mu \FF'$
and $\PP \contr \encoConV {\FF'}$;

\item If  $\FF$ is passive and 
$ \encoConV \FF \arrPI{\mu} \PP$, 
then 
 there is $\FF'$ such that 
$\FF \arrA{\mu} \FF'$
and $\PP  \wbPI \encoConV {\FF'}$.
\end{enumerate} 
\end{lemma} 

\begin{proof}
The proof goes by a case analysis on the possible transition from 
$ \encoConV \FF$, again 
exploiting Lemmas~\ref{l:opt_op}-\ref{l:VAL}.
\end{proof} 
\noindent
We can strengthen both  
Lemma~\ref{l:opcorrCON_str} and 
Theorem~\ref{t:opcorrCON_w}
 to a correspondence between transitions from
\AOGS configurations and from their \piI translation under the (more
restrictive) \opLTS.  
Further, as discussed in Remark~\ref{r:expa_op_ord} and
Appendix~\ref{a:behav}, we can use the expansion relation and
bisimilarity on the
\opLTS, written $\expaOP$ and $\wbOPI$. 
We write \piIop to refer to \piI under the \opLTS.



\begin{theorem}[%
strong transitions, from \AOGS to \piIop]
\label{t:opcorrCON_str_OP}
$ $ \begin{enumerate}

\item If $\FF \longrightarrowA \FF'$, then 
$ \encoConV \FF \raN  \; \: \contrOP  \encoConV {\FF'}$;


\item If $\FF \arrA\ell \FF'$, then
$ \encoConV \FF \arrN\ell  
\; \: \wbOPI 
\encoConV {\FF'}$.
\end{enumerate} 
\end{theorem} 


\begin{theorem} %[from OGS to \piI]
[weak transitions, from \AOGS to \piIop]
\label{t:opcorrCON_w_OP}
$ $ \begin{enumerate}

\item If $\FF \LongrightarrowA \FF'$, then 
$\encoConI \FF \RaN  \; \contrOP \encoConI {\FF'}$;


\item If $\FF \ArrA\ell \FF'$, then
$ \encoConI \FF \ArrN\ell  
\; \wbOPI
\encoConI {\FF'}$.
\end{enumerate} 
\end{theorem} 
\noindent
The following Theorem~\ref{t:opcorrCON3}  is needed in the proof of Theorem~\ref{t:opcorrCON3_w}:
it relates strong transitions from  \piIop processes to
transitions of \AOGS configurations. 
Again, in both theorems, the occurrences of $\expa$ and $\wbPI$ can be
replaced by $\expaOP$ and $\wbOPI$. 

\begin{theorem} %[from \piI to  OGS, sequential LTS]
\label{t:opcorrCON3}
$ $ \begin{enumerate}
% \item  If
% $\encoConV{\FF } \arrN  {\initName (p)} \PP$, then there is $\FF'$ 
% such that
%  $\FF \arr{\initName (p)} \FF'$ 
% and $\PP = \encoConV{\FF' }$ ;

\item If 
$ \encoConV \FF \arrN \tau \PP$ then there is $\FF'$ such that 
$\FF \longrightarrowA \FF'$  and  
$\PP   \contr \encoConV {\FF'}$; 


\item If
$ \encoConV \FF \arrN\ell \PP$
 then there is $\FF'$ such that
 $\FF \arrA\ell \FF'$
and $\PP \wbPI \encoConV {\FF'}$.
\end{enumerate} 
\end{theorem} 
\noindent
We report now more details on the proof of Corollary~\ref{c:traces}.
First, we recall the assertion:

\begin{center}
\begin{tabular}{l}
For any trace $s$, we have 
  $\FF \ArrA s$   iff 
 $ \encoEnvI \FF  \ArrN s$.
\end{tabular} 
\end{center} 


\begin{proof}
The result is first established for 
 $\qencV$, and then extended to 
$\cbvSymb$,  exploiting the correctness of the optimisations 
in  $\qencV$ (Lemma~\ref{l:opt_sound}; again, in this lemma, the occurrence of $\expa$
can be lifted to $\expaOP$).

For  $\qencV$, 
both directions are proved proceeding by induction on the length of a trace 
 $\ell_1, \ldots,   \ell_n $.  
In the direction from left to right, we use 
 Theorem~\ref{t:opcorrCON_w_OP}(2). 
For the converse direction, we 
proceed similarly, this time relying on 
Theorem~\ref{t:opcorrCON3_w}(2) 
 (again, the version with $\wbOPI$ in place of 
 $\wbPI$).
%  $\expa$
% can be lifted to $\expaOP$)
% rely 
% on  Theorem~\ref{t:opcorrCON3}. 
% Consider a 
% process $\encoConV \FF$ and 
% a
% trace $\ell_1 \ldots  \ell_n $ for it.  This means that we have processes $P_1, \ldots,
% P_n$ with 
% $$\encoConV \FF \ArrN{\ell_1} P_1 \ldots \ArrN{\ell_n} P_n $$
% Let $m$ be the number of internal transitions that are thus performed.
% We prove that  also $\FF$ has the trace 
% $\ell_1 ;\ldots ; \ell_n $
%   reasoning by  a lexicographical induction on the pair $(m,n)$
% (setting
%  $(m_1,n_1) < (m_2,n_2)$ if  $n_1< n_2$ or ($ n_1=n_2$ and $m_1<m_2$)).
% In the induction, the appearance of the  expansion relation $\contr$  in clause (2),
% rather than  the bisimilarity  equivalence $\wbPI$, is essential.  
% \finish{this could be expanded, if we want} 
 \end{proof} 




\section{Auxiliary results for Section~\ref{ss:compa_cogs_piI}}
\label{a:aux_tr_bisi_sing}

We report here the auxiliary results needed to establish the 
relationship between Concurrent OGS (\COGS) and \piI. 
%
%
%
%
% The two lemmas below are needed in the proofs of Theorems~\ref{t:COGS_pi} and \ref{t:uptoComp}. 
% The reasoning in their proofs is similar to that for Lemmas~\ref{l:opcorrCON_str} and
% \ref{l:opcorrCON2}.  
% \begin{lemma}[from \COGS to \piI, strong transitions]
% \label{l:opcorrCOGS}
% $ $ \begin{enumerate}
%
% \item If $\FF \longrightarrowp \FF'$, then 
% $ \encoConV \FF \longrightarrowPI  \contr \encoConV {\FF'}$;
%
%
% \item If $\FF \arrC\ell \FF'$, then
% $ \encoConV \FF \arrPI\ell  
% \wbPI
% \encoConV {\FF'}$;
% \end{enumerate} 
% \end{lemma} 
%
% \begin{lemma}[from \piI to  \COGS, strong transitions]
% \label{l:opcorrCOGS2}
% $ $ \begin{enumerate}
%
%
% \item If 
% $ \encoConV \FF \longrightarrowPI \PP$ then there is $\FF'$ such that 
% $\FF \longrightarrowp \FF'$  and  
% $\PP   \contr \encoConV {\FF'}$; 
%
%
% \item If
% $ \encoConV \FF \arrPI\ell \PP$
%  then there is $\FF'$ such that
%  $\FF \arrC\ell \FF'$
% and $\PP \wbPI \encoConV {\FF'}$.
%
% \end{enumerate} 
% \end{lemma} 
 %
%
%
 %           
The results below are used in the proof of Lemma~\ref{l:tr_bisi_sing}. 



A process $\PP$  \emph{has deterministic immediate transitions} if, for any $\act$, 
whenever $\PP \arrPI\act \PP_1 $ and   $\PP \arrPI\act \PP_2$, then $\PP_1 = \PP_2$. 

\begin{lemma}
\label{l:immediate_tr}
Suppose $\PP,\QQ$ have  deterministic immediate transitions, $\PP \TEPI \QQ$, 
 and $\PP \arrPI\act \PP_1$. Then 
$\QQ \ArrPI \act \QQ_1$ implies $\PP _1 \TEPI \QQ_1$. 
\end{lemma} 
\noindent
Lemma below is the analogue to  Lemma~\ref{l:deco} for traces.

\begin{lemma}
\label{l:par_tr}
$ $ 
\begin{itemize}
\item
Suppose $P = P_1 | ! \inp a {\tilb}.P_2$, 
and $Q = Q_1 | ! \inp a {\tilb}.Q_2$, 
and  $\pairP P \TEPI \pairP Q$.
Then also 
  $\pairP {P_1}  \TEPI \pairP {Q_1}$.


\item 
Similarly, suppose $P = P_1 |  \inp a {\tilb}.P_2$, 
and $Q = Q_1 |  \inp a {\tilb}.Q_2$, 
and  $\pairP {P} \TEPI \pairP {Q}$.
Then also 
  $\pairP {P_1}  \TEPI \pairP {Q_1}$.
\end{itemize}
 \end{lemma} 
\noindent
We can now conclude the proof of 
Lemma~\ref{l:tr_bisi_sing}.
We recall the assertion:
\begin{center}
\begin{tabular}{l}
For any $M,N$ we have:  \\
$
{\encoEnvI M} \TEPI {\encoEnvI N}
$ iff 
${\encoEnvI M} \wbPI {\encoEnvI N}$.
\end{tabular} 
\end{center} 

\begin{proof}
We have to prove that trace equivalence
implies bisimilarity. 
We work with the optimised encoding  $\qencV$.
The relation 
\[ 
(\encoConV \FF,  \encoConV \GG) \st \mbox{ $\FF,\GG$ are singleton
  with $\FF \TEp 
\GG$}
\]
is a bisimulation up-to context and up-to $(\contr,
 \wbPI)$. 


 $\FF$ and $\GG$, as singleton, are of the  forms 
described in Lemma~\ref{l:shapesT}.
Moreover, to be in the relation $ \TEp$ they must be of the same
form, and, using also 
  Lemmas~\ref{l:opt_op}  and 
\ref{l:opt_con} we deduce that they have deterministic immediate
transitions. 

Suppose  $\FF$ and $\GG$ are of the form 
$\encoIa {M} p$ and 
$\encoIa {N} q$, respectively.

Consider the case 
$\encoIa {M} p \arr\tau P$; then by Lemma~\ref{l:opt_op}, 
$P \contr \encoIa {M'} p$, for some $M'$. 
Moreover, we have $ \encoIa {M'} p \TEp \encoIa {M} p \TEp  \encoIa
{N} q$, and we are done. 


The case when $\encoIa {M} p$ performs an output action $\bout px$ is
simple, as usual using  Lemma~\ref{l:opt_op}; we also deduce that
$p=q$. 

Suppose now $\encoIa {M} p$ performs an output action $\bout x{z,p'}$.
Using  Lemma~\ref{l:opt_op}, we have 
\[\encoIa {M} p\arr{\bout x{z,p'}}  
\exn \encoVVa \val z|\inp {p'} y.\encoVa {\evctxt[y]} p
\]
for some $\val,y,E$.
Since $\encoIa {M}p \TEp \encoIa {N}q$, and again using 
 Lemma~\ref{l:opt_op},
we deduce 
\[\encoIa {N} q\Longrightarrow \arr{\bout x{z,p'}}  
\exn \encoVVa W z|\inp {q'} y.\encoVa {E'[y]} q
\]
By Lemma~\ref{l:immediate_tr}, 
\[
\begin{array}{rc}
\encoVVa \val z|\inp {p'} y.\encoVa {\evctxt[y]} p
& \TEp \\[3pt]
\encoVVa W z|\inp {q'} y.\encoVa {E'[y]} q
\end{array}
 \] 
Applying Lemma~\ref{l:par_tr} twice, we deduce 
\[
\encoVVa \val z 
\TEp 
\encoVVa W z
 \]
 and
\[
\inp {q'} y.\encoVa {E'[y]} q \TEp
\inp {q'} y.\encoVa {E'[y]} q
\]
and we are done, up to expansion and context. 


\noindent The cases when
$\FF$ and $\GG$ are of a different form are similar, this time using 
Lemma~\ref{l:opt_con}.
\end{proof}              

   


% \section{Auxiliary results for Section~\ref{ss:upto_games}}
% \label{a:ss:upto_games}

% In Definition~\ref{d:uptoComp} 
% we have introduced 
% the technique of  \emph{bisimulation up-to
%  reduction and   composition}.



% Sometimes we do not need  the `closure under reduction'; that is,  referring to  the clauses
% of Definition~\ref{d:uptoComp} 
%  we may take  
% $\FF'_i = \FF_i$, $i=1,2$).
% In this case  
% we say that $\R$ is a 
% `{bisimulation up-to
%    composition}'.
% %  and just take, in clause (1) above   (and similarly for
% % its converse (2)) 
% % $\FF'_1 \uptoComp \R \FF'_2$.





% We now   present an additional up-to technique, which is used in  
% Appendix~\ref{s:enf} when relating 
% OGS bisimilarity  with  eager normal form bisimilarity.
% When  $\R$ is a relation on singleton configuration, a straitforward simplification 
% of `bisimulation up-to
%  reduction and   composition'
% allows
% us to play the bisimulation game only on visible actions; the bisimulation clause becomes: 
% \begin{itemize}
% \item
%  if $\FF_1 \Longrightarrowp\arrC\ell \FF_1'  $ then there is $\FF_2' $ such that 
% $\FF_2\Longrightarrowp\arrC\ell \FF_2'  $ and 
% $\FF'_1 \uptoComp \R \FF'_2$.
% \end{itemize} 
%  We call this a \emph{singleton bisimulation up-to composition}; its soundness is a
%  corollary of that for bisimulation   
%  up-to
%  reduction and  composition, 
%  derived 
% from
%  the up-to techniques for \piI and the
% mapping from \COGS  onto \piI. 






% \begin{proposition}
% \label{p:sing_bis}
% If $\R$ is a singleton bisimulation up-to composition, then ${\R} \subseteq {\wbC}$.
% \end{proposition} 

% \begin{proof}
%  Take 
% $$\S \defi \R \cup \{ (\FF_1,\FF_2) \st 
% \begin{array}[t]{l}
% \mbox{ there are } \FF'_1, \FF'_2  \mbox{ with   } \FF'_1 \RR \FF'_2 
%  \\
% \mbox{ and   }\FF'_i \Longrightarrowp \FF_i  \mbox{ $i=1,2$}  \}
% \end{array} 
% $$
% Intuitively,  $\S$  is the closure of $\R$ under deterministic reductions. 
% Then $\S$ is a bisimulation up-to composition. 
% (Note: we exploit the fact that transitions of singleton configurations are deterministic.)
% \end{proof} 



% %%%%%END MOVED 

% We now  report  more details on the proof of Theorem~\ref{t:uptoComp}.
% We recall its assertion: 

% \myteo{t:uptoComp}
% If 
% $\R$ is bisimulation up-to
%  reduction and  composition then ${\R} \subseteq {\wbC}$.

% We first need the following
% Lemma~\ref{l:uptoOGSpi}, which is a  consequence of the compositionality of the mapping from
% \COGS to \piI.  
% \begin{lemma}
% \label{l:uptoOGSpi}
% For a relation $\R$ on \COGS configurations, 
% if 
% ${(\FF_1, \FF_2)} \in  { \uptoComp \R}$, then 
% we have $( \encoConV {\FF_1},  \encoConV {\FF_2}) \in  \ctx{( \encoConV{ \R })}$. 
% \end{lemma} 

% We are now ready to show  the proof of Theorem~\ref{t:uptoComp}.
% \begin{proof}
% We prove the result via \piI, exploiting Lemma~\ref{l:ogsp_pi_traces_bis}(2).
% Here again, it is convenient to use
% the  
% optimised encoding 
% $\qencV$.

% If $\R$ is a relation on \COGS configurations, then 
% $\encoConV \R$ is the relation on \piI processes obtained by mapping each pair
% in $\R$ in the expected manner: 
% \[ \encoConV \R \defi \{ (\encoConV {\FF_1}, \encoConV {\FF_2}) \st \FF_1 \RR \FF_2 
%  \} \] 
% And, if $\S$ is a relation on \piI processes, then $\ctx \S$ is the closure of $\S$ under
% polyadic contexts, i.e., the set of pairs  of the form  
% $(\ct {P_1,\ldots, P_n},\ct {Q_1,\ldots, Q_n}$ where $\qct$ is a polyadic context and for
% each $i$, $P_i \SS Q_i$. 

% As a   consequence of the compositionality of the mapping from
% \COGS to \piI,
% for  a relation   $\R$ on \COGS configurations, 
% if 
% ${(\FF_1, \FF_2)} \in  { \uptoComp \R}$, then 
%  $( \encoConV {\FF_1},  \encoConV {\FF_2}) \in  \ctx{( \encoConV{ \R })}$. 




% Now, let $\R$ be the relation in the hypothesis of the theorem.
% We  prove the theorem by  showing that 
% $\encoConV \R$ is a bisimulation
%  up-to context and up-to $(\contr,
%  \wbPI)$ in 
% \piI, and then appealing to Theorem~\ref{t:up-to-con-expa-nointer}.  

% Suppose 
% $\encoConV {\FF_1}  \arrPI \ell P$.  Then
% there is $\FF_1'$ with  
% $\FF_1  \arrPI \ell \FF_1' $ and $P  \wbPI \encoConV {\FF'_1}$ (this is derived from Lemma~\ref{l:opcorrCOGS}). 
% Since $\R$ is bisimulation up-to
%  reduction and  composition and $\FF_1 \RR \FF_2$, there are 
%   $\FF_1''$, and $\FF_2'$, with  
% $\FF'_1 \Longrightarrowp \FF''_1$, 
% $\FF_2
%   \arrPI \ell \Longrightarrowp \FF_2' $ and
% $\FF''_1 \uptoComp \R \FF'_2$. 


% Using  Theorem~\ref{t:COGS_pi}(1) 
% %Lemma~\ref{l:Redp} 
% from 
% $\FF'_1 \Longrightarrowp \FF''_1$ and the inclusion ${\contr} \subseteq
% {\wbPI}$, 
% we derive
% $\encoConV {\FF'_1}  \wbPI  \encoConV {\FF''_1}$; 
% and  by   
% Theorem~\ref{t:COGS_pi}(1),
% %Lemmas~\ref{l:Redp}, 
% and  Lemmas \ref{l:opt_sound} and \ref{l:opcorrCOGS},  from 
% $\FF_2
%   \arrPI \ell \Longrightarrowp \FF_2' $ 
% we derive
% $\encoConV {\FF_2}  \arrPI \ell \wbPI  \encoConV {\FF'_2}$.   


% Finally,  
% \iflong
% by Lemma~\ref{l:uptoOGSpi},  
% \fi
% from $(\FF''_1, \FF'_2) \in  \uptoComp \R$, we derive 
%  $( \encoConV {\FF'_1},  \encoConV {\FF'_2}) \in  \ctx{( \encoConV{ \R })}$. 
% This closes the proof,  up-to polyadic contexts and bisimilarity.

% The case when 
% $\encoConV {\FF_1}$ makes a silent move is  simpler,  using expansion in place of
% bisimilarity.
% \end{proof} 


% %%%%%END MOVED 
